\chapter{Introduction}

\section{Domain: sugarcane farming in Karnataka}
Sugarcane is one of the main commercial crops of Karnataka. \todo{Add area, production
and number of sugarcane farmers in Karnataka with a citation, e.g. from the Directorate of
Sugarcane Development status paper or the Karnataka agriculture department. Do not use a
number without a source.}
It is a long-duration crop (about 10 to 18 months in the field), and the farmer
makes many decisions over that time: which variety to plant for the season and region,
how to treat setts before planting, how much fertiliser to apply and when, how to manage
weeds, and how to identify and control pests and diseases such as early shoot borer,
termites, smut and red rot. Recommendations differ by region: the University of
Agricultural Sciences (UAS) Bengaluru publishes a Package of Practices for southern
Karnataka and UAS Dharwad publishes one for northern Karnataka, and the two do not
always agree.

\section{Problems faced by farmers}
\begin{itemize}
  \item \textbf{Access to expert advice.} The official recommendations exist, but in long
    printed books (in the UAS Bengaluru 2025 Package of Practices, sugarcane is chapter 34,
    starting after 170 pages of other crops), and extension officers are few compared with the number of farmers.
    \todo{citation for extension worker to farmer ratio, if used}
  \item \textbf{Language.} Most farmers in Karnataka ask questions in Kannada. Much of the
    digital agricultural content and most AI tools work best in English.
  \item \textbf{Wrong information is costly.} Advice from shop owners, neighbours or the
    internet may name a product that is not registered for sugarcane, a banned
    pesticide, or a dose for a different pest. A wrong pesticide or dose can destroy the
    crop, harm the person spraying it, or leave residues.
  \item \textbf{Regional differences.} A dose that is right for one region or season can be
    wrong for another.
\end{itemize}

\section{How AI can help farmers}
Natural language processing lets a farmer ask a question in their own words, in their own
language, and get an answer at any time on a mobile phone, instead of searching a printed
book or waiting for a call centre. Retrieval-augmented generation (RAG) lets a language
model answer from a chosen set of documents instead of only from what it learned during
training. If the documents are the official university recommendations, the answers can
in principle be grounded in trusted sources.

\section{Existing approaches and their drawbacks}
\subsection{Methods}
\begin{itemize}
  \item \textbf{Telephone helplines.} The Kisan Call Centre (1800-180-1551) gives
    farmers access to experts in their language, but depends on call-centre capacity and
    working hours.
  \item \textbf{Rule-based and FAQ chatbots.} Answer a fixed list of questions with
    fixed answers. Safe, but they fail on any wording they were not built for.
  \item \textbf{General-purpose LLM chatbots.} Fluent in many languages and answer almost
    anything, but they can state wrong facts with confidence and do not say where an
    answer came from.
  \item \textbf{RAG-based agricultural assistants.} Retrieve passages from an agricultural
    document collection and ask an LLM to answer from them.
    \todo{Cite 3--5 agricultural chatbot / RAG papers from the literature survey here.}
\end{itemize}

\subsection{Drawbacks}
\begin{itemize}
  \item An LLM can still change a number, combine two passages, or answer a question about
    a pest the documents do not cover by using the closest passage it finds. We observed
    exactly this in our own version~1 (Chapter~\ref{ch:results}).
  \item Retrieval returns the \emph{most similar} passage, not a \emph{correct} one: a
    question about an unknown pest still retrieves the nearest pest's dose.
  \item Common text-overlap metrics (ROUGE, BLEU, chrF) used to evaluate such systems do not
    detect a wrong number, and ROUGE-L does not work on Kannada text
    (Section~\ref{sec:metric-failures}).
  \item The source documents themselves are often of unknown origin or mix regions. The
    corpus this project started from contained Tamil Nadu recommendations and a termite
    dose that actually belonged to early shoot borer.
\end{itemize}

\section{Problem statement}
To build a Kannada question-answering assistant for sugarcane farmers in Karnataka that
gives answers only from official, region-specific recommendations, shows where each
answer comes from, and refuses (pointing the farmer to the Kisan Call Centre) whenever it
cannot give a verified answer, so that it never gives a wrong pesticide, dose or unit.

\section{Objectives}
\begin{enumerate}
  \item Build a Kannada chat interface and backend that answers sugarcane cultivation
    questions.
  \item Build and evaluate a RAG pipeline (version~1) for Kannada, and measure where and
    why it fails.
  \item Find evaluation metrics that actually work for Kannada answers and for safety, and
    show which commonly used metrics do not.
  \item Build an answer path (version~2) that gives doses and chemicals only from records
    confirmed by a person against the printed source page, and refuses everything else.
  \item Evaluate both versions on held-out questions, including questions that must be
    refused.
\end{enumerate}

\section{Scope of the study}
\begin{itemize}
  \item \textbf{Crop:} sugarcane only.
  \item \textbf{Region:} Karnataka; southern recommendations from UAS Bengaluru (2025),
    northern from UAS Dharwad (2020). The two are kept separate and never merged.
  \item \textbf{Topics:} varieties, planting season, seed rate, sett treatment,
    fertiliser, irrigation, weed control, major pests and diseases, harvest and ratoon.
  \item \textbf{Out of scope:} market prices, weather, government schemes, other crops,
    diagnosis from symptoms or photos, and voice input. The assistant refuses these
    questions.
\end{itemize}

\section{Usefulness}
A farmer can ask in Kannada and get the university recommendation for their region,
with the book and page it came from, or a clear statement that the assistant does not
know, together with the Kisan Call Centre number. The project also gives a reusable
lesson for low-resource-language agricultural assistants: how to measure safety honestly,
and why a refusal is a correct output.
